\section{Unimodular heights: proof of Theorem~\ref{thm:U}}\label{app:unimodular}

Theorem~\ref{thm:U} asks, for every admissible $\Delta$, for a unimodular triangulation $\mathcal T$ of $\partial\Delta^\circ$
refining its faces and an integral height $\check h$ on the fan over $\mathcal T$ such that $\check h$ and
$\check h-\check\varphi$ are strictly convex. Section~\ref{ssec:ureform} shows that this is equivalent to the existence of
a fine, regular, unimodular triangulation of $\Delta^\circ$ that is star-shaped from the origin, and that strict convexity
of $\check h-\check\varphi$ implies that of $\check h$. Section~\ref{ssec:ufacets} explains why admissibility does not
yield such a triangulation by a general argument: it constrains the facets of $\Delta^\circ$ only along their edges, and
lattice polytopes of dimension three need not have unimodular triangulations at all. Section~\ref{ssec:uproof} proves the
theorem by exhaustion over the classification of Kreuzer and Skarke, with a certificate for each polytope checked in
exact integer arithmetic, and Section~\ref{ssec:utrust} states on what the proof rests.

Throughout, a piecewise linear function $g$ on a complete simplicial fan is \emph{strictly convex} if it is convex and its
maximal domains of linearity are the maximal cones. A \emph{wall} is a common facet of two maximal cones $\kappa,\kappa'$,
and the \emph{bend} of $g$ across it is $g(m')-\ell_\kappa(m')$, where $\ell_\kappa$ is the linear function equal to $g$ on
$\kappa$ and $m'$ is the generator of $\kappa'$ not in $\kappa$. Bends are additive in $g$, and $g$ is strictly convex if
and only if every bend is positive; this is the wall criterion of \cite[\S6.1]{CLS}, where support functions carry the
opposite sign, so that convexity here is concavity there. Recall that $\check\varphi$ equals $1$ on $\partial\Delta^\circ$ and is
linear on the cones over the faces of $\Delta^\circ$; it is convex, its maximal domains of linearity are the cones over the
facets, and $\check\varphi(m)=\max_v\langle m,-v\rangle$ over the vertices $v$ of $\Delta$.

\subsection{Reformulation}\label{ssec:ureform}

\begin{lemma}\label{lem:Ureform}
Let $\mathcal T$ be a unimodular triangulation of $\partial\Delta^\circ$ refining its faces. The following are
equivalent.
\begin{enumerate}
\item[(a)] There is a piecewise linear function $\check h$ on the fan over $\mathcal T$, integral on $M$, such that $\check h$ and
$\check h-\check\varphi$ are strictly convex.
\item[(b)] There is a real piecewise linear function on the fan over $\mathcal T$ that is strictly convex.
\item[(c)] There are real heights $\omega$ on the vertices of $\mathcal T$ such that for every facet $F$ of $\Delta^\circ$ the
restriction $\mathcal T|_F$ is the regular subdivision of $\operatorname{vert}(\mathcal T)\cap F$ induced by $\omega|_F$.
\end{enumerate}
Moreover, for any piecewise linear $\check h$ on the fan over $\mathcal T$, strict convexity of $\check h-\check\varphi$
implies strict convexity of $\check h$.
\end{lemma}

\begin{proof}
(a)$\Rightarrow$(b): take $\check h-\check\varphi$.

(b)$\Rightarrow$(a): strict convexity is an open condition on the values at the rays, so we may take the strictly convex
function $g$ rational, and after scaling integral at the primitive generators of the rays. Put $\check h=g+\check\varphi$.
Then $\check h-\check\varphi=g$ is strictly convex. Every simplex of $\mathcal T$ lies in a facet, so $\check\varphi$ is
linear on every cone of the fan and bends by $\ge0$ across every wall; hence every bend of $\check h$ is at least the
corresponding bend of $g$, and $\check h$ is strictly convex. The same computation proves the last assertion, for any triangulation. The
function $g$ is integral on $M$, because a linear function on a unimodular cone that is integral at the generators is
integral on the lattice points of the cone, and $\check\varphi$ is integral on $M$ because the vertices of $\Delta$ are
lattice points.

(b)$\Rightarrow$(c): put $\omega=g|_{\operatorname{vert}\mathcal T}$. On the cone over a facet $F$, $g$ is convex with the cones
over the simplices of $\mathcal T|_F$ as its maximal domains of linearity. Its restriction to the affine hyperplane of $F$
is therefore a convex piecewise linear function on $F$ whose domains of linearity are exactly the simplices of
$\mathcal T|_F$ and whose values at their vertices are $\omega$; this says that $\mathcal T|_F$ is the regular subdivision
induced by $\omega|_F$.

(c)$\Rightarrow$(b): let $g_0$ be linear on each cone of the fan over $\mathcal T$ with $g_0=\omega$ on the vertices. Across a
wall inside the cone over a facet $F$, the bend of $g_0$ is positive by the regularity of $\mathcal T|_F$, since a bend is
a local condition, homogeneous of degree one. Across the finitely many walls between the cones over two different facets,
$\check\varphi$ bends positively, while it does not bend across walls inside a facet cone. So $g_0+\lambda\check\varphi$
is strictly convex for $\lambda\gg0$.
\end{proof}

Adding a constant to the heights $\omega$ on $\partial\Delta^\circ$ adds a multiple of $\check\varphi$ to the corresponding
piecewise linear function; this is the conversion used below. A triangulation of $\Delta^\circ$ is \emph{star-shaped from
the origin} if the origin is a vertex of each of its maximal simplices.

\begin{corollary}\label{cor:Ureform}
The conclusion of Theorem~\ref{thm:U} holds for $\Delta$ if and only if $\Delta^\circ$ has a fine, regular, unimodular
triangulation that is star-shaped from the origin.
\end{corollary}

\begin{proof}
A star-shaped triangulation of $\Delta^\circ$ consists of the simplices $\operatorname{conv}(0,\tau)$, $\tau$ in a
triangulation $\mathcal T$ of $\partial\Delta^\circ$; each $\tau$ lies in a facet, since a convex subset of the boundary of a
polytope lies in a facet, so $\mathcal T$ refines the faces of $\Delta^\circ$. Since $\Delta^\circ$ is reflexive, every facet lies at lattice distance
one from the origin, so $\operatorname{conv}(0,\tau)$ is unimodular if and only if $\tau$ is unimodular in its facet, and a
unimodular simplex contains no lattice points other than its vertices, so unimodular triangulations are fine. The
star-shaped triangulation is regular if and only if it is induced by heights $0$ at the origin and $\omega$ on
$\partial\Delta^\circ$ (subtract an affine function), and the lower envelope of the lifted points is then linear on each
$\operatorname{conv}(0,\tau)$ and $0$ at the origin; so regularity holds if and only if the piecewise linear function with
values $\omega$ on the fan over $\mathcal T$ is strictly convex. This is (b) of Lemma~\ref{lem:Ureform}, which is equivalent to (a).
\end{proof}

\subsection{The facets of \texorpdfstring{$\Delta^\circ$}{the polar polytope}}\label{ssec:ufacets}
Let $v$ be a vertex of $\Delta$ with neighbours $w_1,\dots,w_k$ along edges, and $d_i=w_i-v$, primitive by admissibility.
The facet of $\Delta^\circ$ dual to $v$ is the lattice three-polytope
\[
F_v=\{m\in M_\R:\ \langle m,v\rangle=-1,\ \langle m,d_i\rangle\ge0\ (i=1,\dots,k)\},
\]
the polytope of the anticanonical bundle of $\mathbb P_\Delta$ restricted to the toric divisor $D_v$. Its vertices, edges
and two-faces correspond to the facets, two-faces and edges of $\Delta$ that contain $v$.

\begin{proposition}\label{prop:Ufacets}
Let $\Delta$ be admissible and let $E$ be an edge of $F_v$, dual to a two-face $\theta\ni v$ of $\Delta$.
\begin{enumerate}
\item The normal cone of $F_v$ along $E$ is lattice equivalent to the tangent cone of $\theta$ at $v$, in the lattice of the
plane of $\theta$. It is unimodular if $\theta$ is a triangle, a square or a hexagon, and if $\theta$ is a pentagon and $v$
is one of its three corners of index one; at the other two corners of a pentagon it has index $2$.
\item If $\theta$ is not a triangle, $E$ has lattice length one.
\end{enumerate}
\end{proposition}

\begin{proof}
(1) The two-dimensional cone of the fan $\operatorname{Star}(v)/v$ given by $\theta$ is $\operatorname{cone}(\theta-v)$ modulo
$\R v$. Since $\theta$ lies in a facet of $\Delta$ at lattice height one, the saturated lattice of $\R\theta$ is
$\Z v\oplus L_\theta$, with $L_\theta$ the direction lattice of the plane of $\theta$, so $L_\theta$ maps isomorphically onto a
saturated sublattice of $N/\Z v$, carrying the tangent cone of $\theta$ at $v$ to that two-dimensional cone. The index at a
corner is the determinant of the two primitive edge vectors there; for the pentagon
$\operatorname{conv}\{(1,0),(1,1),(0,1),(-1,0),(0,-1)\}$ these are $1,1,1,2,2$, for the hexagon all are $1$. (2) is the
condition $\ell(\theta^\circ)=1$ of Definition~\ref{def:admissible}.
\end{proof}

Admissibility is a condition on the two-faces of $\Delta$ and on their dual edges, that is, on the edges of the facets
$F_v$ and on the normal cones along them; Proposition~\ref{prop:Ufacets} records what it says there. It imposes no direct
condition on the cones of $F_v$ at its vertices beyond their two-dimensional faces, nor on the two-faces of $F_v$.

The known sufficient conditions for unimodular triangulations do not apply to these facets. If every facet $F_v$ has
lattice width one with respect to each of its facet normals (is \emph{compressed}), then every pulling triangulation
of each facet is unimodular \cite[Thm.~2.3]{HPPS}, and pulling the boundary lattice points of $\Delta^\circ$ in one global
order gives heights as in Lemma~\ref{lem:Ureform}(c). If the facet normals of a facet form a unimodular matrix, that facet
has a regular unimodular triangulation \cite[Thm.~2.4]{HPPS}, but the triangulations of two adjacent facets obtained in this
way need not agree on the common face. Compressed facets are rare: of the $3{,}175$ facets of the $418$ admissible
polytopes with at most eight vertices, $46$ are compressed. Facets with unimodular normals are common ($12{,}350$ of the
$13{,}469$ facets of the admissible polytopes of the list of Section~\ref{ssec:list}), but their triangulations need not be
compatible. The facets are also far from small: among the $3{,}175$ there are facets with $242$ lattice points and
normalised volume $756$, and $2{,}171$ of the $2{,}444$ facets with at most $40$ lattice points contain empty lattice
tetrahedra of volume greater than one among their lattice points, so that a fine triangulation of such a facet need not be
unimodular. For lattice
three-polytopes in general, unimodular triangulations need not exist, and their existence is open even for smooth
lattice three-polytopes \cite[\S1.5.1]{HPPS}; the facets $F_v$ need not be smooth, since by
Proposition~\ref{prop:Ufacets} their vertex cones are constrained only through their two-dimensional faces. A general
proof of Theorem~\ref{thm:U} would therefore produce compatible regular unimodular triangulations for a class of lattice
three-polytopes, the polytopes of $-K_{\mathbb P_\Delta}|_{D_v}$, for which no such result is known. We know no conceptual
proof, and we prove the theorem by exhaustion.

\subsection{Proof of Theorem~\ref{thm:U}}\label{ssec:uproof}
The proof has three steps: the list of all reflexive four-polytopes, the determination of the admissible ones, and an
exactly checked certificate for each of these.

\emph{Step 1: the classification.} Kreuzer and Skarke classified the reflexive four-polytopes up to $GL_4(\Z)$: there are
$473{,}800{,}776$ of them \cite{KS00}. We use a copy of their list whose completeness we checked against the published
count, with every polytope reflexive and no two equivalent (by normal forms). Every polytope of the list is treated as
$\Delta$; since the list contains the polar of each of its members, this covers every polar pair in both orders.

\emph{Step 2: the admissible polytopes.} Admissibility is decided from the vertices and the facet inequalities of
$\Delta$, computed with PALP \cite{PALP}. The two-faces of $\Delta$ are the maximal vertex sets shared by two facets; each is
classified in the lattice of its plane. By Pick's theorem, a lattice polygon with primitive edges is a unimodular
triangle, a unit square, a $dP_7$ pentagon or a $dP_6$ hexagon exactly when its number of vertices and twice its area are
$(3,1)$, $(4,2)$, $(5,5)$ or $(6,6)$. For every non-triangular two-face the lattice length of the dual edge is computed
from the two facet normals. The result is Table~\ref{tab:Ucounts}: up to $GL_4(\Z)$ there are $51{,}827$ admissible
polytopes, $40{,}403$ without and $11{,}424$ with a hexagonal two-face, and $741$ of them have no singular two-face. Every one of them was checked admissible a second time by an independent
program, the $51{,}827$ have pairwise different normal forms, and the $612$ admissible polytopes of the list of
Section~\ref{ssec:list} all occur among them.

\begin{table}[ht]
\centering\small\setlength{\tabcolsep}{4pt}
\begin{tabular}{rrrr@{\qquad}rrrr}
vertices & polytopes & admissible & with hexagon & vertices & polytopes & admissible & with hexagon\\\hline
5 & 1{,}561 & 5 & 0 & 20 & 21{,}913{,}680 & 3{,}407 & 1{,}557\\
6 & 24{,}189 & 26 & 0 & 21 & 12{,}070{,}919 & 2{,}404 & 1{,}268\\
7 & 177{,}446 & 102 & 0 & 22 & 5{,}826{,}221 & 1{,}567 & 906\\
8 & 834{,}638 & 285 & 1 & 23 & 2{,}450{,}720 & 999 & 623\\
9 & 2{,}867{,}955 & 629 & 7 & 24 & 898{,}929 & 583 & 386\\
10 & 7{,}725{,}801 & 1{,}156 & 25 & 25 & 284{,}696 & 315 & 221\\
11 & 16{,}608{,}387 & 1{,}931 & 57 & 26 & 78{,}468 & 173 & 131\\
12 & 29{,}270{,}253 & 2{,}857 & 123 & 27 & 18{,}417 & 81 & 65\\
13 & 43{,}458{,}000 & 3{,}888 & 225 & 28 & 3{,}781 & 38 & 31\\
14 & 56{,}060{,}584 & 4{,}846 & 345 & 29 & 647 & 14 & 12\\
15 & 64{,}085{,}869 & 5{,}499 & 527 & 30 & 114 & 9 & 7\\
16 & 65{,}615{,}931 & 5{,}824 & 782 & 31 & 23 & 4 & 4\\
17 & 59{,}972{,}682 & 5{,}697 & 1{,}088 & 32 & 8 & 0 & 0\\
18 & 48{,}703{,}033 & 5{,}134 & 1{,}419 & 33 & 2 & 1 & 1\\
19 & 34{,}847{,}821 & 4{,}352 & 1{,}612 & 36 & 1 & 1 & 1\\\hline
& & & & total & $473{,}800{,}776$ & $51{,}827$ & $11{,}424$
\end{tabular}
\caption{Reflexive four-polytopes by number of vertices, the admissible ones, and those among them with a hexagonal
two-face.}\label{tab:Ucounts}
\end{table}

\emph{Step 3: certificates.} For each admissible $\Delta$ we record an integer height $h(m)$ at every nonzero lattice point
$m$ of $\Delta^\circ$. It determines a candidate triangulation: in each facet $F$ of
$\Delta^\circ$, the lower faces of the convex hull of the lifted points $(m,h(m))$, $m\in F\cap M$. For an integer $\lambda$, computed in the check, put $g=h+\lambda$
on $\partial\Delta^\circ\cap M$, extended linearly on the cones over the candidate tetrahedra, and $\check h=g+\check\varphi$.
The following are then checked in exact integer arithmetic:
\begin{enumerate}
\item[(i)] every recorded point is a lattice point of $\partial\Delta^\circ$, and every height is an integer;
\item[(ii)] every candidate tetrahedron lies in a facet hyperplane of $\Delta^\circ$ and has determinant $\pm1$;
\item[(iii)] every triangle of a candidate tetrahedron lies in exactly two candidate tetrahedra, and their fourth vertices
lie on opposite sides of the hyperplane spanned by the triangle and the origin;
\item[(iv)] each of three fixed rays in general position lies in the cone over exactly one candidate tetrahedron;
\item[(v)] across every wall, the bends of $g$ and of $\check h$ are positive integers.
\end{enumerate}

\begin{lemma}\label{lem:Ucert}
If (i)--(v) hold, the candidate tetrahedra form a unimodular triangulation $\mathcal T$ of $\partial\Delta^\circ$ refining its
faces, and $\check h$ is an integral height on the fan over $\mathcal T$ with $\check h$ and $\check h-\check\varphi$ strictly
convex.
\end{lemma}

\begin{proof}
Write $\sigma_T=\operatorname{cone}(T)$ for a candidate tetrahedron $T$; by (ii) it is a four-dimensional simplicial cone, and
its facets are the cones over the triangles of $T$, the \emph{walls}. Let $Z$ be the union of the cones over the edges of
the candidate tetrahedra and of the intersections of pairs of walls that span different hyperplanes; it is a finite union of
cones of dimension at most two, so $M_\R\setminus Z$ is connected, and any two points off $Z$ and off the walls are joined by
a path that avoids $Z$ and crosses the walls transversally at finitely many points.

\emph{The interiors are disjoint.} For $x$ off the walls let $N(x)$ be the number of cones $\sigma_T$ containing $x$. Let $x$
lie on a wall but not in $Z$. The walls through $x$ contain $x$ in their relative interiors (the relative boundary of a
wall lies in cones over edges) and span a common hyperplane $H$. A cone containing $x$ on its boundary contains it in
exactly one facet (two facets meet in the cone over an edge), which is one of these walls; by (iii) each of these walls is
a facet of exactly two cones, lying on opposite sides of $H$. So near $x$ the same number of cones is gained on each side of
$H$, and $N$ takes the same value on both sides. Hence $N$ is constant off the walls, and by (iv) it equals $1$: the
interiors of the cones are pairwise disjoint.

\emph{Local structure at the faces.} We show, for a face $\phi\neq0$ of some $\sigma_T$ and a point $x$ in its relative interior,
that the cones having $\phi$ as a face cover a neighbourhood of $x$, and that no other cone contains $x$. The second
assertion follows from the first: a cone containing $x$ has interior points arbitrarily close to $x$, which would lie in the
interior of a cone having $\phi$ as a face, so by disjointness it is that cone. If $\phi=\sigma_T$ there is nothing to show.
If $\phi$ is a wall, by (iii) it is a facet of a second cone on the other side of its hyperplane, and the two cones cover a
neighbourhood of $x$; in particular a wall is never subdivided by the facets of several cones on one side, since (iii)
requires the whole wall to be a facet of a cone on each side. If $\phi$ is the cone over an edge $e$, project to the plane
$M_\R/\operatorname{span}\phi$: the cones containing $\phi$ become sectors of angle less than $\pi$, and the walls
containing $\phi$ their boundary rays. By (iii) each such ray bounds exactly two of the sectors, on opposite sides, so the
sectors form cycles that turn monotonically around the origin, each covering every direction at least once; their
interiors are disjoint, so there is exactly one cycle and it covers every direction once. So the cones containing $\phi$
cover a neighbourhood of $x$. If $\phi$ is a ray $\R_{\ge0}v$, project to $M_\R/\R v$ and intersect with a small sphere: the
cones containing $\phi$ become spherical triangles forming a closed two-dimensional pseudomanifold $L$ by (iii). By the case
of walls and edges already treated, the map from $L$ to the sphere is a local homeomorphism near every point of $L$,
including its vertices, whose links are single circles. So $L$ is a closed surface and the map is a covering of the
two-sphere; each component of $L$ maps homeomorphically, and by disjointness of the interiors there is exactly one. So the
cones containing $\phi$ cover a neighbourhood of $x$.

\emph{The fan.} The union of the cones is closed, and by the local structure it is open, so it is $M_\R$. Let
$x\in\sigma_T\cap\sigma_{T'}$, $x\neq0$, and let $\phi$ be the face of $\sigma_T$ containing $x$ in its relative interior; by
the local structure $\phi$ is a face of $\sigma_{T'}$. Applying this to a point $x$ in the relative interior of the convex
cone $\sigma_T\cap\sigma_{T'}$, whose carrier face $\phi$ in $\sigma_T$ contains $\sigma_T\cap\sigma_{T'}$, gives
$\sigma_T\cap\sigma_{T'}=\phi$, a common face. So the cones form a complete simplicial fan.

\emph{Conclusion.} By (i) and (ii) each tetrahedron lies in a facet of $\Delta^\circ$ and is unimodular, so the tetrahedra
triangulate every facet, hence every face, of $\partial\Delta^\circ$; a unimodular tetrahedron contains no lattice points
besides its vertices, and the tetrahedra cover $\partial\Delta^\circ$, so every lattice point of $\partial\Delta^\circ$ is a
vertex. By (v) and the wall criterion, $g=\check h-\check\varphi$ and $\check h$ are strictly convex. The values of $g$ at the
vertices are integers, so $\check h$ is integral on $M$ as in the proof of Lemma~\ref{lem:Ureform}.
\end{proof}

All $51{,}827$ certificates pass, with $12{,}001{,}172$ tetrahedra in all. By Lemma~\ref{lem:Ucert} each
certificate proves the conclusion of Theorem~\ref{thm:U} for its polytope, and by Step 2, which examined every polytope of
the classification, the certified polytopes are all admissible reflexive four-polytopes. This proves
Theorem~\ref{thm:U}. \qed

A second program proves the conclusion from the same heights by a different set of sufficient checks, in exact rational
arithmetic: (i$'$) the recorded points are exactly the nonzero lattice points of $\Delta^\circ$, enumerated independently,
and lie on its boundary; (ii$'$) in each facet $F$, every cell of the lower convex hull of the points $(m,h(m))$,
$m\in F\cap M$, is a tetrahedron of determinant $\pm1$; (v$'$) for each tetrahedron and each recorded point $m$ outside its
cone, $g(m)$ and $\check h(m)$ exceed the values at $m$ of the linear functions that agree with $g$, respectively $\check h$,
on the cone. By (ii$'$) each facet is triangulated by the regular subdivision induced by $h$, and since the heights are
global, the triangulations of two facets restrict to the same regular subdivision of their common face; so the tetrahedra
form a unimodular triangulation of $\partial\Delta^\circ$ refining its faces, and (v$'$) is strict convexity of $g$ and
$\check h$ on its fan. This program does not need (iii) and (iv). As a further consistency check it compares the number of
tetrahedra with the normalised volume of $\Delta^\circ$.

The heights were found by choosing $h(m)$ close to a large multiple of a positive definite quadratic form in $m$, first
$|m|^2$ and then random forms: the first choice succeeds on $7{,}231$ polytopes, and no polytope needed more than $30$
forms. A quadratic height induces a regular triangulation of every facet at once, which is why one height serves all
facets; the constant $\lambda$ makes the bends across walls between different facets positive, as in the proof of
Lemma~\ref{lem:Ureform}. How the heights were found plays no role in the proof.

\subsection{What the proof rests on}\label{ssec:utrust}
The proof rests on three computations.
\begin{enumerate}
\item \emph{The classification data.} The theorem is as reliable as the classification of \cite{KS00}, itself a computer
classification, and as our copy of it. The classification was not recomputed. Our copy was checked against the published
count and published file digests; every member was checked reflexive, and no two members are equivalent.
\item \emph{The admissibility filter.} The program of Step 2 was run on every polytope of the classification, and the
statement that the certified polytopes are all the admissible ones rests on it. Every polytope it accepts was checked
admissible again by an independent program. A second implementation of the whole test, sharing no code with the first
(its faces come from a different polyhedral library, and the interior lattice points of two-faces from Pick's theorem), was
run on part of the classification and agrees with the first on every polytope on which both were run. Let $d$ be the
number of nonzero lattice points of $\Delta$ that are not vertices; all admissible polytopes found have $d\le7$ except two,
with $d=12$ and $d=20$.
The second implementation has been run on every polytope with at most $9$ vertices ($3{,}905{,}789$ polytopes), and, for
$10$ to $13$ vertices, on every polytope with $d\le5$ together with $5{,}000$ further polytopes for each number of vertices
($2{,}000$ chosen uniformly and $3{,}000$ with $d\in\{6,7\}$): $7{,}644{,}706$ polytopes, containing all $10{,}879$ admissible
polytopes with at most $13$ vertices. For $14$ or more vertices, which carry $40{,}948$ of the $51{,}827$ admissible
polytopes, every polytope accepted by the first program was checked admissible by the second, and the rejected
polytopes rest on the first program alone.
\item \emph{The certificate checks.} The first program checks (i)--(v) in integer arithmetic, with determinants by
fraction-free elimination; the second checks (i$'$), (ii$'$) and (v$'$) in exact rational arithmetic. Both accept all
$51{,}827$ certificates. On $135$ corrupted certificates (random, negated, constant or rescaled heights, a point removed)
the two programs return the same verdicts.
\end{enumerate}
The certificates and both checking programs are available at \coderepo.
